% Original pages: 121-128
% Figures: none; three mathematical diagrams typeset directly
% Uncertainty log: none

iii) There exists $n_0$ such that $M_{n_0+r}=\mathfrak a^rM_{n_0}$ for all $r\geq0$, hence $G(M)$ is
generated by $\bigoplus_{n\leq n_0}G_n(M)$.  Each $G_n(M)=M_n/M_{n+1}$ is Noetherian and
annihilated by $\mathfrak a$, hence is a finitely-generated $A/\mathfrak a$-module, hence $\bigoplus_{n\leq n_0}G_n(M)$
is generated by a finite number of elements (as an $A/\mathfrak a$-module), hence $G(M)$ is
finitely generated as a $G(A)$-module.\proofbox

The last main result of this chapter is that the $\mathfrak a$-adic completion of a
Noetherian ring is Noetherian.  Before we can proceed to the proof we need a
simple lemma connecting the completion of any filtered group and the associated graded group.

\medskip
\noindent\textit{\textbf{Lemma 10.23.\amtarget{10.23}} Let $\phi:A\longrightarrow B$ be a homomorphism of filtered groups, i.e.
$\phi(A_n)\subseteq B_n$, and let $G(\phi):G(A)\longrightarrow G(B)$, $\widehat\phi:\widehat A\longrightarrow\widehat B$ be the induced homomorphisms of the associated graded and completed groups.  Then}

\noindent\textit{i) $G(\phi)$ injective $\Rightarrow\widehat\phi$ injective;}

\noindent\textit{ii) $G(\phi)$ surjective $\Rightarrow\widehat\phi$ surjective.}

\noindent\emph{Proof.}  Consider the commutative diagram of exact sequences
\[
\begin{array}{ccccccccc}
0&\longrightarrow&A_n/A_{n+1}&\longrightarrow&A/A_{n+1}&\longrightarrow&A/A_n&\longrightarrow&0\\
&&\big\downarrow{\scriptstyle G_n(\phi)}&&\big\downarrow{\scriptstyle\alpha_{n+1}}&&\big\downarrow{\scriptstyle\alpha_n}&&\\
0&\longrightarrow&B_n/B_{n+1}&\longrightarrow&B/B_{n+1}&\longrightarrow&B/B_n&\longrightarrow&0.
\end{array}
\]
This gives the exact sequence
\[
\begin{split}
0\longrightarrow\operatorname{Ker}G_n(\phi)\longrightarrow\operatorname{Ker}\alpha_{n+1}\longrightarrow\operatorname{Ker}\alpha_n
\longrightarrow\operatorname{Coker}G_n(\phi)\longrightarrow\operatorname{Coker}\alpha_{n+1}\\
\longrightarrow\operatorname{Coker}\alpha_n\longrightarrow0.
\end{split}
\]
From this we see, by induction on $n$, that $\operatorname{Ker}\alpha_n=0$ (case i)) or $\operatorname{Coker}\alpha_n=0$
(case ii)).  Moreover in case ii) we also have $\operatorname{Ker}\alpha_{n+1}\longrightarrow\operatorname{Ker}\alpha_n$ surjective.
Taking the inverse limit of the homomorphisms $\alpha_n$ and applying \amref{10.2} the lemma
follows.\proofbox

We can now form a result which is a partial converse of \amref{10.22} iii) and is the
main step in showing that $\widehat A$ is Noetherian.

\medskip
\noindent\textit{\textbf{Proposition 10.24.\amtarget{10.24}} Let $A$ be a ring, $\mathfrak a$ an ideal of $A$, $M$ an $A$-module, $(M_n)$ an
$\mathfrak a$-filtration of $M$.  Suppose that $A$ is complete in the $\mathfrak a$-topology and that $M$ is
Hausdorff in its filtration topology (i.e. that $\bigcap_n M_n=0$).  Suppose also that
$G(M)$ is a finitely-generated $G(A)$-module.  Then $M$ is a finitely-generated
$A$-module.}

\noindent\emph{Proof.}  Pick a finite set of generators of $G(M)$, and split them up into their
homogeneous components, say $\xi_i$ $(1\leq i\leq v)$ where $\xi_i$ has degree say $n(i)$, and is
therefore the image of say $x_i\in M_{n(i)}$.  Let $F^i$ be the module $A$ with the stable
$\mathfrak a$-filtration given by $F_k^i=\mathfrak a^{k+n(i)}$ and put $F=\bigoplus_{i=1}^vF^i$.  Then mapping the
generator 1 of each $F^i$ to $x_i$ defines a homomorphism
\[
\phi:F\longrightarrow M
\]
of filtered groups, and $G(\phi):G(F)\longrightarrow G(M)$ is a homomorphism of $G(A)$
modules.  By construction it is surjective.  Hence by \amref{10.23} ii) $\widehat\phi$ is surjective.
Consider now the diagram
\[
\begin{array}{ccc}
F&\xrightarrow{\phi}&M\\
{\scriptstyle\alpha}\big\downarrow&&\big\downarrow{\scriptstyle\beta}\\
\widehat F&\xrightarrow{\widehat\phi}&\widehat M
\end{array}
\]
Since $F$ is free and $A=\widehat A$ it follows that $\alpha$ is an isomorphism.  Since $M$ is
Hausdorff $\beta$ is injective.  The surjectivity of $\widehat\phi$ thus implies the surjectivity of $\phi$,
and this means that $x_1,\ldots,x_v$ generate $M$ as an $A$-module.\proofbox

\medskip
\noindent\textit{\textbf{Corollary 10.25.\amtarget{10.25}} With the hypotheses of \amref{10.24}, if $G(M)$ is a Noetherian
$G(A)$-module, then $M$ is a Noetherian $A$-module.}

\noindent\emph{Proof.}  We have to show that every submodule $M'$ of $M$ is finitely generated
\amref{6.2}.  Let $M'_n=M'\cap M_n$; then $(M'_n)$ is an $\mathfrak a$-filtration of $M'$, and the embedding $M'_n\longrightarrow M_n$ gives rise to an injective homomorphism $M'_n/M'_{n+1}\longrightarrow
M_n/M_{n+1}$, hence to an embedding of $G(M')$ in $G(M)$.  Since $G(M)$ is Noetherian,
$G(M')$ is finitely generated by \amref{6.2}; also $M'$ is Hausdorff, since $\bigcap M'_n\subseteq
\bigcap M_n=0$; hence by \amref{10.24} $M'$ is finitely generated.\proofbox

We can now deduce the result we are after:

\medskip
\noindent\textit{\textbf{Theorem 10.26.\amtarget{10.26}} If $A$ is a Noetherian ring, $\mathfrak a$ an ideal of $A$, then the $\mathfrak a$-completion $\widehat A$ of $A$ is Noetherian.}

\noindent\emph{Proof.}  By \amref{10.22} we know that
\[
G_\mathfrak a(A)=G_{\widehat{\mathfrak a}}(\widehat A)
\]
is Noetherian.  Now apply \amref{10.25} to the complete ring $\widehat A$, taking $M=\widehat A$
(filtered by $\widehat{\mathfrak a}^n$, and so Hausdorff).\proofbox

\medskip
\noindent\textit{\textbf{Corollary 10.27.\amtarget{10.27}} If $A$ is a Noetherian ring, the power series ring $B=
A[[x_1,\ldots,x_n]]$ in $n$ variables is Noetherian.  In particular $k[[x_1,\ldots,x_n]]$
($k$ a field) is Noetherian.}

\noindent\emph{Proof.}  $A[x_1,\ldots,x_n]$ is Noetherian by the Hilbert basis theorem, and $B$ is its
completion for the $(x_1,\ldots,x_n)$-adic topology.\proofbox

\subsection{Exercises}

\noindent\textbf{1.} Let $\alpha_n:\mathbb Z/p\mathbb Z\longrightarrow\mathbb Z/p^n\mathbb Z$ be the injection of abelian groups given by $\alpha_n(1)=p^{n-1}$,
and let $\alpha:A\longrightarrow B$ be the direct sum of all the $\alpha_n$ (where $A$ is a countable direct
sum of copies of $\mathbb Z/p\mathbb Z$, and $B$ is the direct sum of the $\mathbb Z/p^n\mathbb Z$).  Show that the
$p$-adic completion of $A$ is just $A$ but that the completion of $A$ for the topology
induced from the $p$-adic topology on $B$ is the direct product of the $\mathbb Z/p\mathbb Z$.  Deduce
that $p$-adic completion is not a right-exact functor on the category of all $\mathbb Z$-modules.

\medskip
\noindent\textbf{2.} In Exercise 1, let $A_n=\alpha^{-1}(p^nB)$, and consider the exact sequence
\[
0\longrightarrow A_n\longrightarrow A\longrightarrow A/A_n\longrightarrow0.
\]
Show that $\varprojlim$ is not right exact, and compute $\varprojlim{}^1A_n$.

\medskip
\noindent\textbf{3.} Let $A$ be a Noetherian ring, $\mathfrak a$ an ideal and $M$ a finitely-generated $A$-module.
Using Krull's Theorem and Exercise 14 of Chapter 3, prove that
\[
\bigcap_{n=1}^\infty\mathfrak a^nM=\bigcap_{\mathfrak m\supseteq\mathfrak a}\operatorname{Ker}(M\longrightarrow M_\mathfrak m),
\]
where $\mathfrak m$ runs over all maximal ideals containing $\mathfrak a$.

Deduce that
\[
\widehat M=0\quad\Longleftrightarrow\quad\operatorname{Supp}(M)\cap V(\mathfrak a)=\varnothing\qquad(\text{in }\operatorname{Spec}(A)).
\]
[The reader should think of $\widehat M$ as the ``Taylor expansion'' of $M$ transversal to
the subscheme $V(\mathfrak a)$: the above result then shows that $M$ is determined in a
neighborhood of $V(\mathfrak a)$ by its Taylor expansion.]

\medskip
\noindent\textbf{4.} Let $A$ be a Noetherian ring, $\mathfrak a$ an ideal in $A$, and $\widehat A$ the $\mathfrak a$-adic completion.  For
any $x\in A$, let $\widehat x$ be the image of $x$ in $\widehat A$.  Show that
\[
x\text{ not a zero-divisor in }A\quad\Rightarrow\quad\widehat x\text{ not a zero-divisor in }\widehat A.
\]
Does this imply that
\[
A\text{ is an integral domain}\quad\Rightarrow\quad\widehat A\text{ is an integral domain}?
\]
[Apply the exactness of completion to the sequence $0\longrightarrow A\xrightarrow{x}A$.]

\medskip
\noindent\textbf{5.} Let $A$ be a Noetherian ring and let $\mathfrak a,\mathfrak b$ be ideals in $A$.  If $M$ is any $A$-module, let
$M^\mathfrak a,M^\mathfrak b$ denote its $\mathfrak a$-adic and $\mathfrak b$-adic completions respectively.  If $M$ is finitely
generated, prove that $(M^\mathfrak a)^\mathfrak b\cong M^{\mathfrak a+\mathfrak b}$.

[Take the $\mathfrak a$-adic completion of the exact sequence
\[
0\longrightarrow\mathfrak b^mM\longrightarrow M/\mathfrak b^mM\longrightarrow0
\]
and apply \amref{10.13}.  Then use the isomorphism
\[
\varprojlim_m\left(\varprojlim_n M/(\mathfrak a^nM+\mathfrak b^mM)\right)\cong
\varprojlim_n M/(\mathfrak a^nM+\mathfrak b^nM)
\]
and the inclusions $(\mathfrak a+\mathfrak b)^{2n}\subseteq\mathfrak a^n+\mathfrak b^n\subseteq(\mathfrak a+\mathfrak b)^n$.]

\medskip
\noindent\textbf{6.} Let $A$ be a Noetherian ring and $\mathfrak a$ an ideal in $A$.  Prove that $\mathfrak a$ is contained in the
Jacobson radical of $A$ if and only if every maximal ideal of $A$ is closed for the
$\mathfrak a$-topology.  (A Noetherian topological ring in which the topology is defined by
an ideal contained in the Jacobson radical is called a \emph{Zariski ring}.  Examples are
local rings and (by \amref{10.15}(iv)) $\mathfrak a$-adic completions.)

\medskip
\noindent\textbf{7.} Let $A$ be a Noetherian ring, $\mathfrak a$ an ideal of $A$, and $\widehat A$ the $\mathfrak a$-adic completion.  Prove
that $\widehat A$ is faithfully flat over $A$ (Chapter 3, Exercise 16) if and only if $A$ is a Zariski
ring (for the $\mathfrak a$-topology).

[Since $\widehat A$ is flat over $A$, it is enough to show that
\[
M\longrightarrow\widehat M\text{ injective for all finitely generated }M\quad\Longleftrightarrow\quad A\text{ is Zariski};
\]
now use \amref{10.19} and Exercise 6.]

\medskip
\noindent\textbf{8.} Let $A$ be the local ring of the origin in $\mathbb C^n$ (i.e., the ring of all rational functions
$f/g\in\mathbb C(z_1,\ldots,z_n)$ with $g(0)\ne0$), let $B$ be the ring of power series in $z_1,\ldots,z_n$
which converge in some neighborhood of the origin, and let $C$ be the ring of
formal power series in $z_1,\ldots,z_n$, so that $A\subset B\subset C$.  Show that $B$ is a local
ring and that its completion for the maximal ideal topology is $C$.  Assuming that
$B$ is Noetherian, prove that $B$ is $A$-flat.  [Use Chapter 3, Exercise 17, and
Exercise 7 above.]

\medskip
\noindent\textbf{9.} Let $A$ be a local ring, $\mathfrak m$ its maximal ideal.  Assume that $A$ is $\mathfrak m$-adically complete.
For any polynomial $f(x)\in A[x]$, let $\bar f(x)\in(A/\mathfrak m)[x]$ denote its reduction mod. $\mathfrak m$.
Prove \emph{Hensel's lemma}: if $f(x)$ is monic of degree $n$ and if there exist coprime
monic polynomials $\bar g(x),\bar h(x)\in(A/\mathfrak m)[x]$ of degrees $r,n-r$ with $\bar f(x)=
\bar g(x)\bar h(x)$, then we can lift $\bar g(x),\bar h(x)$ back to monic polynomials $g(x),h(x)\in A[x]$
such that $f(x)=g(x)h(x)$.

[Assume inductively that we have constructed $g_k(x),h_k(x)\in A[x]$ such that
$g_k(x)h_k(x)-f(x)\in\mathfrak m^kA[x]$.  Then use the fact that since $\bar g(x)$ and $\bar h(x)$ are
coprime we can find $\bar a_p(x),\bar b_p(x)$, of degrees $\leq n-r,r$ respectively, such that
$x^p=\bar a_p(x)\bar g_k(x)+\bar b_p(x)\bar h_k(x)$, where $p$ is any integer such that $1\leq p\leq n$.
Finally, use the completeness of $A$ to show that the sequences $g_k(x),h_k(x)$
converge to the required $g(x),h(x)$.]

\medskip
\noindent\textbf{10.} i) With the notation of Exercise 9, deduce from Hensel's lemma that if $\bar f(x)$
has a simple root $\alpha\in A/\mathfrak m$, then $f(x)$ has a simple root $a\in A$ such that
$\alpha=a\bmod\mathfrak m$.

ii) Show that 2 is a square in the ring of 7-adic integers.

iii) Let $f(x,y)\in k[x,y]$, where $k$ is a field, and assume that $f(0,y)$ has $y=a_0$
as a simple root.  Prove that there exists a formal power series $y(x)=
\sum_{n=0}^\infty a_nx^n$ such that $f(x,y(x))=0$.

(This gives the ``analytic branch'' of the curve $f=0$ through the point $(0,a_0)$.)

\medskip
\noindent\textbf{11.} Show that the converse of \amref{10.26} is false, even if we assume that $A$ is local and
that $\widehat A$ is a finitely-generated $A$-module.

[Take $A$ to be the ring of germs of $C^\infty$ functions of $x$ at $x=0$, and use Borel's
Theorem that every power series occurs as the Taylor expansion of some $C^\infty$
function.]

\medskip
\noindent\textbf{12.} If $A$ is Noetherian, then $A[[x_1,\ldots,x_n]]$ is a faithfully flat $A$-algebra.  [Express
$A\longrightarrow A[[x_1,\ldots,x_n]]$ as a composition of flat extensions, and use Exercise 5(v)
of Chapter 1.]

\section{Dimension Theory}

One of the basic notions in algebraic geometry is that of the dimension of a
variety.  This is essentially a local notion, and, as we shall show in this chapter,
there is a very satisfactory theory of dimension for general Noetherian local
rings.  The main theorem asserts the equivalence of three different definitions of
dimension.  Two of these definitions have a fairly obvious geometrical content,
but the third involving the Hilbert function is less conceptual.  It has, however,
many technical advantages and the whole theory becomes more streamlined if
one brings it in at an early stage.

After dealing with dimension we give a brief account of regular local rings,
which correspond to the notion of non-singularity in algebraic geometry.  We
establish the equivalence of three definitions of regularity.

Finally we indicate how, in the case of algebraic varieties over a field, the
local dimensions we have defined coincide with the transcendence degree of the
function field.

\subsection{Hilbert functions}

Let $A=\bigoplus_{n=0}^\infty A_n$ be a Noetherian graded ring.  By \amref{10.7} $A_0$ is a Noetherian
ring, and $A$ is generated (as an $A_0$-algebra) by say $x_1,\ldots,x_s$, which we may take
to be homogeneous, of degrees $k_1,\ldots,k_s$ (all $>0$).

Let $M$ be a finitely-generated graded $A$-module.  Then $M$ is generated by a
finite number of homogeneous elements, say $m_j$ $(1\leq j\leq t)$; let $r_j=\deg m_j$.
Every element of $M_n$, the homogeneous component of $M$ of degree $n$, is thus of
the form $\sum_j f_j(x)m_j$, where $f_j(x)\in A$ is homogeneous of degree $n-r_j$ (and
therefore zero if $n<r_j$).  It follows that $M_n$ is finitely generated as an $A_0$-module, namely it is generated by all $g_j(x)m_j$ where $g_j(x)$ is a monomial in the
$x_i$ of total degree $n-r_j$.

Let $\lambda$ be an \emph{additive function} (with values in $\mathbb Z$) on the class of all finitely-generated $A_0$-modules (Chapter 2).  The \emph{Poincar\'e series} of $M$ (with respect to $\lambda$)
is the generating function of $\lambda(M_n)$, i.e., it is the power series
\[
P(M,t)=\sum_{n=0}^\infty\lambda(M_n)t^n\in\mathbb Z[[t]].
\]

\medskip
\noindent\textit{\textbf{Theorem 11.1.\amtarget{11.1}} (Hilbert, Serre). $P(M,t)$ is a rational function in $t$ of the form
$f(t)/\prod_{i=1}^s(1-t^{k_i})$, where $f(t)\in\mathbb Z[t]$.}

\noindent\emph{Proof.}  By induction on $s$, the number of generators of $A$ over $A_0$.  Start with
$s=0$; this means that $A_n=0$ for all $n>0$, so that $A=A_0$ and $M$ is a finitely-generated $A_0$ module, hence $M_n=0$ for all large $n$.  Thus $P(M,t)$ is a polynomial
in this case.

Now suppose $s>0$ and the theorem true for $s-1$.  Multiplication by $x_s$
is an $A$-module homomorphism of $M_n$ into $M_{n+k_s}$, hence it gives an exact
sequence, say
\begin{equation}
0\longrightarrow K_n\longrightarrow M_n\xrightarrow{x_s}M_{n+k_s}\longrightarrow L_{n+k_s}\longrightarrow0.
\tag{1}
\end{equation}
Let $K=\bigoplus_nK_n$, $L=\bigoplus_nL_n$; these are both finitely-generated $A$-modules
(because $K$ is a submodule and $L$ a quotient module of $M$), and both are annihilated by $x_s$, hence they are $A_0[x_1,\ldots,x_{s-1}]$-modules.  Applying $\lambda$ to (1) we
have, by \amref{2.11}
\[
\lambda(K_n)-\lambda(M_n)+\lambda(M_{n+k_s})-\lambda(L_{n+k_s})=0;
\]
multiplying by $t^{n+k_s}$ and summing with respect to $n$ we get
\begin{equation}
(1-t^{k_s})P(M,t)=P(L,t)-t^{k_s}P(K,t)+g(t)
\tag{2}
\end{equation}
where $g(t)$ is a polynomial.  Applying the inductive hypothesis the result now
follows.\proofbox

The order of the pole of $P(M,t)$ at $t=1$ we shall denote by $d(M)$.  It
provides a measure of the ``size'' of $M$ (relative to $A$).  In particular $d(A)$ is
defined.  The case when all $k_i=1$ is specially simple:

\medskip
\noindent\textit{\textbf{Corollary 11.2.\amtarget{11.2}} If each $k_i=1$, then for all sufficiently large $n$, $\lambda(M_n)$ is a
polynomial in $n$ (with rational coefficients) of degree\footnote{We adopt the convention here that the degree of the zero polynomial is $-1$: also
that the binomial coefficient $\binom n{-1}=0$ for $n\geq0$, and $=1$ for $n=-1$.} $d-1$.}

\noindent\emph{Proof.}  By \amref{11.1} we have $\lambda(M_n)={}$ coefficient of $t^n$ in $f(t)\cdot(1-t)^{-s}$.  Canceling powers of $(1-t)$ we may assume $s=d$ and $f(1)\ne0$.  Suppose $f(t)=
\sum_{k=0}^N a_kt^k$; since
\[
(1-t)^{-d}=\sum_{k=0}^\infty\binom{d+k-1}{d-1}t^k
\]
we have
\[
\lambda(M_n)=\sum_{k=0}^Na_k\binom{d+n-k-1}{d-1}\quad\text{for all }n\geq N.
\]
and the sum on the right-hand side is a polynomial in $n$ with leading term
$(\sum a_k)n^{d-1}/(d-1)!\ne0$.\proofbox

\emph{Remarks.}  1) For a polynomial $f(x)$ to be such that $f(n)$ is an integer for all
integers $n$, it is not necessary for $f$ to have integer coefficients: e.g., $\frac12x(x+1)$.

2) The polynomial in \amref{11.2} is usually called the \emph{Hilbert function} (or polynomial) of $M$ (with respect to $\lambda$).

Returning now to the sequence (1) let us replace $x_s$ by any element $x\in A_k$
which is not a zero-divisor in $M$ (i.e., $xm=0$ with $m\in M\Rightarrow m=0$).  Then
$K=0$ and equation (2) shows that
\[
d(L)=d(M)-1.
\]
Thus we have proved

\medskip
\noindent\textit{\textbf{Proposition 11.3.\amtarget{11.3}} If $x\in A_k$ is not a zero-divisor in $M$ then $d(M/xM)=
d(M)-1$.}\proofbox

We shall use \amref{11.1} in the case where $A_0$ is an Artin ring (in particular, a
field) and $\lambda(M)$ is the \emph{length} $l(M)$ of a finitely-generated $A_0$-module $M$.  By
\amref{6.9} $l(M)$ is additive.

\noindent\textbf{Example.}  Let $A=A_0[x_1,\ldots,x_s]$, where $A_0$ is an Artin ring and the $x_i$ are
independent indeterminates.  Then $A_n$ is a free $A_0$-module generated by the
monomials $x_1^{m_1}\cdots x_s^{m_s}$ where $\sum m_i=n$; there are $\binom{s+n-1}{s-1}$ of these, hence
$P(A,t)=(1-t)^{-s}$.

We shall now consider the Hilbert functions obtained from a local ring by
passing to the associated graded rings as in Chapter 10.

\medskip
\noindent\textit{\textbf{Proposition 11.4.\amtarget{11.4}} Let $A$ be a Noetherian local ring, $\mathfrak m$ its maximal ideal,
$\mathfrak q$ an $\mathfrak m$-primary ideal, $M$ a finitely-generated $A$-module, $(M_n)$ a stable $\mathfrak q$-filtration of $M$.  Then}

\noindent\textit{i) $M/M_n$ is of finite length, for each $n\geq0$;}

\noindent\textit{ii) for all sufficiently large $n$ this length is a polynomial $g(n)$ of degree $\leq s$
in $n$, where $s$ is the least number of generators of $\mathfrak q$;}

\noindent\textit{iii) the degree and leading coefficient of $g(n)$ depend only on $M$ and $\mathfrak q$, not on
the filtration chosen.}

\noindent\emph{Proof.}  i) Let $G(A)=\bigoplus_n\mathfrak q^n/\mathfrak q^{n+1}$, $G(M)=\bigoplus_nM_n/M_{n+1}$.  $G_0(A)=A/\mathfrak q$ is an
Artin local ring, say by \amref{8.5}; $G(A)$ is Noetherian, and $G(M)$ is a finitely-generated graded $G(A)$-module \amref{10.22}.  Each $G_n(M)=M_n/M_{n+1}$ is a
Noetherian $A$-module annihilated by $\mathfrak q$, hence a Noetherian $A/\mathfrak q$-module, and
therefore of finite length (since $A/\mathfrak q$ is Artin).  Hence $M/M_n$ is of finite length,
and
\begin{equation}
l_n=l(M/M_n)=\sum_{r=1}^n l(M_{r-1}/M_r).
\tag{1}
\end{equation}

ii) If $x_1,\ldots,x_s$ generate $\mathfrak q$, the images $\bar x_i$ of the $x_i$ in $\mathfrak q/\mathfrak q^2$ generate $G(A)$
as an $A/\mathfrak q$-algebra, and each $\bar x_i$ has degree 1.  Hence by \amref{11.2} we have $l(M_n/M_{n+1})
=f(n)$ say, where $f(n)$ is a polynomial in $n$ of degree $\leq s-1$ for all large $n$.

Since from (1) we have $l_{n+1}-l_n=f(n)$, it follows that $l_n$ is a polynomial $g(n)$
of degree $\leq s$, for all large $n$.

iii) Let $(\widetilde M_n)$ be another stable $\mathfrak q$-filtration of $M$, and let $\widetilde g(n)=l(M/\widetilde M_n)$.
By \amref{10.6} the two filtrations have bounded difference, i.e., there exists an integer
$n_0$ such that $M_{n+n_0}\subseteq\widetilde M_n$, $\widetilde M_{n+n_0}\subseteq M_n$ for all $n\geq0$; consequently we have
$g(n+n_0)\geq\widetilde g(n)$, $\widetilde g(n+n_0)\geq g(n)$.  Since $g$ and $\widetilde g$ are polynomials for all
large $n$, we have $\lim_{n\to\infty}g(n)/\widetilde g(n)=1$, and therefore $g,\widetilde g$ have the same degree
and leading coefficient.\proofbox

The polynomial $g(n)$ corresponding to the filtration $(\mathfrak q^nM)$ is denoted by
$\chi_\mathfrak q^M(n)$:
\[
\chi_\mathfrak q^M(n)=l(M/\mathfrak q^nM)\qquad(\text{for all large }n).
\]
If $M=A$, we write $\chi_\mathfrak q(n)$ for $\chi_\mathfrak q^A(n)$ and call it the \emph{characteristic polynomial}
of the $\mathfrak m$-primary ideal $\mathfrak q$.  In this case \amref{11.4} gives

\medskip
\noindent\textit{\textbf{Corollary 11.5.\amtarget{11.5}} For all large $n$, the length $l(A/\mathfrak q^n)$ is a polynomial $\chi_\mathfrak q(n)$
of degree $\leq s$, where $s$ is the least number of generators of $\mathfrak q$.}\proofbox

The polynomials $\chi_\mathfrak q(n)$ for different choices of the $\mathfrak m$-primary ideal $\mathfrak q$ all have
the same degree, as the next proposition shows:

\medskip
\noindent\textit{\textbf{Proposition 11.6.\amtarget{11.6}} If $A,\mathfrak m,\mathfrak q$ are as above}
\[
\deg\chi_\mathfrak q(n)=\deg\chi_\mathfrak m(n).
\]
\noindent\emph{Proof.}  We have $\mathfrak m\supseteq\mathfrak q\supseteq\mathfrak m^r$ for some $r$ by \amref{7.16}, hence $\mathfrak m^n\supseteq\mathfrak q^n\supseteq\mathfrak m^{rn}$ and
therefore
\[
\chi_\mathfrak m(n)\leq\chi_\mathfrak q(n)\leq\chi_\mathfrak m(rn)\quad\text{for all large }n.
\]
Now let $n\longrightarrow\infty$, remembering that the $\chi$'s are polynomials in $n$.\proofbox

The common degree of the $\chi_\mathfrak q(n)$ will be denoted by $d(A)$: in view of \amref{11.2}
this means we are putting $d(A)=d(G_\mathfrak m(A))$ where $d(G_\mathfrak m(A))$ is the integer
defined earlier as the pole at $t=1$ of the Hilbert function of $G_\mathfrak m(A)$.

\subsection{Dimension theory of Noetherian local rings}

Let $A$ be a Noetherian local ring, $\mathfrak m$ its maximal ideal.

Let $\delta(A)={}$ least number of generators of an $\mathfrak m$-primary ideal of $A$.

Our ambition is to prove that $\delta(A)=d(A)=\dim A$.  We shall achieve this by
proving $\delta(A)\geq d(A)\geq\dim A\geq\delta(A)$.  \amref{11.5} and \amref{11.6} together provide the
first link in this chain:

\medskip
\noindent\textit{\textbf{Proposition 11.7.\amtarget{11.7}} $\delta(A)\geq d(A)$.}\proofbox

Next we shall prove the analogue for local rings of \amref{11.3}.  Note that this
proof uses the strong version of the Artin--Rees lemma (not just the topological
part).
